\chapter{Tips And Tricks}
\thispagestyle{empty}

\newcommand{\examplecommand}[1]{\textbackslash\texttt{#1}}
\newcommand{\examplespace}[1]{\hspace{2cm}}

\section{Bibliography and citing}
\paragraph*{Bibliography:} For managing your bibliography, we recommend to use programs like:
\begin{itemize}
  \item \texttt{jabref}
  \item Mendeley
\end{itemize}
Use \examplecommand{cite}\{BibTexLabel\} when you refer to literature, like \cite{Bear1988,Flemisch2011,Dreybrodt1996}.
\paragraph*{Citing:} Check your given reference thoroughly! This does not only
mean text references or quotes, but also personal communication, input for
discussion/hypotheses/conclusions, ... .\\
Figures should also be emphasized
here: give references/sources to figures if you have copied or redraw them. Best would
be to even ask for the authors allowance.

\section{General Tips for Writing}
\begin{itemize}
  \item for non-running indices, like w for water are non-italic, use \examplecommand{textrm}\{w\}
        which produces $S_\textrm{w}$ instead of $S_w$.
  \item use the \examplecommand{DuMuX} command to get DuMuX written that way: \DuMuX
  \item for tables and figures use \examplecommand{caption}[ShortText]\{LongText\}
        then ShortText will appear in the list of figures/tables
        and LongText will appear directly below the figure/table
  \item fractions inside text look nicer with the \examplecommand{nicefrac}\{a\}\{b\}
        command \nicefrac{a}{b}
\end{itemize}

\section{Tables, Figures, and Equations}
This section provides some basic examples for tables, figures, and equations.
Table~\ref{tab:exampleTable} shows an exemplary table. 

\begin{table}
% table caption is above the table
\caption{This is an exemplary caption for your table. Table captions are above the actual table.}
\label{tab:exampleTable}       % Give a unique label
\centering
\begin{tabular}{lll}
\hline\noalign{\smallskip}
first & second & third  \\ % these are the column headers
\noalign{\smallskip}\hline\noalign{\smallskip}
number & number & number \\
number & number & number \\
\noalign{\smallskip}\hline
\end{tabular}
\end{table}

Figure~\ref{fig:exampleFigure} is an exemplary figure. 
Always reference your tables and figures at least once in the text. 
For normal figures, a width of approximately 0.8-0.9 of the text width \examplecommand{textwidth} is a good choice.
Be sure that the font size in plots is readable, which means at least the size of \examplecommand{tiny}. 
Every good plot should have at least axis labels and legend and may have a title if necessary.
Make sure that the lines, markers, and colors in your plots are easily visible, consistent through different figures,
meaningful, and distinguishable even for black and white print.

\begin{figure}
\centering
%for .eps, .png, etc use:
\includegraphics[width=0.2\textwidth]{\Pics/Logo_LH2_sw}
%for tikz figures use:
%\input{your_file.tex}
\caption{This is an exemplary caption for your figure. Figure captions are below the actual figure.}
\label{fig:exampleFigure}       % Give a unique label
\end{figure}

Equation~\eqref{eq:massBalance} is an example of a mass balance equation including components
for a multi-phase flow case.
Introduce every new variable in the vicinity of its first appearance. 
Be consistent in the format of the indices, sub- and superscripts and always label equations.

\begin{equation}\label{eq:massBalance} % Give a unique label
 \sum_{\alpha} \left[\frac{\partial}{\partial t}\left(\phi \rho_\alpha x^\kappa_\alpha S_\alpha \right) 
 + \nabla\cdotp \left(\rho_\alpha x^\kappa_\alpha \mathbf{v}_\alpha \right) 
 - \nabla\cdotp \left(\rho_\alpha \mathbf{D}_\mathrm{pm,\alpha} \nabla x^\kappa_\alpha \right) \right]= q^\kappa
 %\:\alpha \in \mathrm{n,w}
\end{equation}  


\section{Further Information}
\begin{enumerate}
  \item A thesis guideline is available at the homepage of Faculty 2:\\
        \href{http://www.f02.uni-stuttgart.de/studium/}{www.f02.uni-stuttgart.de/studium/}
  \item Good answer to many questions regarding LaTeX can be found in the web.
\end{enumerate}

\section{Additional Packages}
\subsection{Units}
This requires the \texttt{units} package. It is helpful to get the right
format for number, text and spacing. It works in both, math and text mode:\\
\examplecommand{unit}[1]\{m\} \hfill $\unit[1]{m}$\\
\examplecommand{unitfrac}[1]\{m\}\{s\} \hfill \unitfrac[1]{m}{s}\\

\subsection{Todonote}
This requires the \texttt{todonote}
\todo{With this package nice todos can be created}{package.}\\
You can even use them inline:
\todo[inline]{This is an inline todo note}


\endinput
